% ============================================================================
% Fichier  : partie3/P03_C16_deontologie_profession.tex
% Titre    : Déontologie et Profession de l'Auditeur
% Auteur   : MINKA MI NGUIDJOÏ Thierry Emmanuel
% Version  : 1.0 -- novembre 2026
% ============================================================================
% Chapitre propre à INF3226 (pas d'équivalent dans ACSI).
% Sources : IIA, Normes mondiales d'audit interne (2024), domaine II
% (éthique et professionnalisme : intégrité, objectivité, compétence,
% diligence professionnelle, confidentialité) ; ISACA, Code of Professional
% Ethics ; IESBA, Code international de déontologie des professionnels
% comptables (cinq principes fondamentaux et cinq menaces) ; ISACA,
% exigences de formation continue des certifiés (CPE). Les règles nationales
% (ONECCA, textes sur l'expertise judiciaire) sont présentées sans citation,
% à vérifier.
% Corrigés : corriges/Corr_C16.tex.
% ============================================================================

\chapter{Déontologie et Profession de l'Auditeur}
\label{chap:deontologie}
\chaptermark{Déontologie et profession}

\epigraph{%
    « La compétence fait un bon auditeur.\\
    La déontologie fait un auditeur qu'on peut croire. »%
}{Principe de la profession}

\niveauChapitre{Fondamental}{%
    Maîtrise des Chapitres~\ref{chap:normes} (normes, indépendance) et
    \ref{chap:droit} (responsabilité, secret)%
}

\begin{boxobjectifs}
\begin{itemize}
    \item Énoncer les principes déontologiques communs aux codes de l'IIA,
          de l'ISACA et de l'IESBA.
    \item Reconnaître les cinq menaces sur l'indépendance et choisir une
          sauvegarde adaptée.
    \item Résoudre un dilemme déontologique par une méthode écrite.
    \item Distinguer les responsabilités civile, pénale et disciplinaire de
          l'auditeur.
    \item Situer les certifications professionnelles de l'audit des SI et
          leurs obligations de formation continue.
    \item Décrire le marché de l'audit et les risques qu'il fait peser sur
          l'indépendance.
\end{itemize}
\end{boxobjectifs}

\bigskip

% ============================================================================
\section*{Question directrice : que faire quand le client a raison de vouloir autre chose~?}
% ============================================================================

Le 12~juin~2026, M.~ETEME, directeur général de LiZbiZ, reçoit la note
d'AUDIX CONSEIL sur les ajustements de stock de Douala et les paiements de
TRANSLOG SERVICES. Il appelle M.~BEYALA~: « Votre note me met dans une
position impossible. Le conseil se réunit dans dix jours, nous sommes en
pleine levée de fonds. Retirez ces deux points du rapport, je m'en occupe en
interne. Et, bien sûr, nous avons d'autres missions pour vous à
l'automne. »

M.~ETEME n'a pas tort de vouloir protéger son entreprise. Mais M.~BEYALA ne
peut pas accepter. Ce chapitre explique pourquoi, et comment refuser sans
rompre.

\medskip
\textbf{Plan du chapitre.}
\begin{enumerate}
    \item Les codes de la profession (section~\ref{sec:c16-codes}).
    \item L'indépendance et ses menaces (section~\ref{sec:c16-menaces}).
    \item Résoudre un dilemme (section~\ref{sec:c16-dilemme}).
    \item La responsabilité de l'auditeur (section~\ref{sec:c16-responsabilite}).
    \item Les certifications et la compétence (section~\ref{sec:c16-certif}).
    \item Le marché de l'audit (section~\ref{sec:c16-marche}).
\end{enumerate}

\separateur

% ============================================================================
\section{Les Codes de la Profession}
\label{sec:c16-codes}
% ============================================================================

\subsection{Trois codes, des principes communs}

Trois codes encadrent la conduite des auditeurs des SI, selon leur
appartenance.

\begin{tableaucours}{L{2.8cm}L{3.6cm}L{5.8cm}}{%
    \tableauHeader{Code} &
    \tableauHeader{Qui y est tenu} &
    \tableauHeader{Principes (reformulés)}
}
IIA, Normes mondiales, domaine~II & Auditeurs internes, membres et certifiés
de l'IIA & Intégrité, objectivité, compétence, diligence
professionnelle, confidentialité~\cite{iia2024gias} \\
ISACA, code de déontologie & Certifiés et membres de l'ISACA (CISA, CISM,
CRISC\ldots) & Soutenir les normes, agir avec diligence et intégrité,
servir l'intérêt des parties, préserver la confidentialité, maintenir
sa compétence, informer des résultats, promouvoir la
profession~\cite{isaca2024ethique} \\
IESBA, code international & Professionnels comptables, dont les
commissaires aux comptes & Intégrité, objectivité, compétence et
diligence, confidentialité, comportement professionnel~\cite{iesba2018code} \\
\end{tableaucours}
\vspace{-0.8em}
\begin{center}{\small\itshape Tableau~16.1. Trois codes de déontologie}\end{center}

Les trois codes convergent sur cinq exigences~: être honnête (intégrité),
juger sans parti pris (objectivité), savoir faire ce qu'on accepte
(compétence), travailler avec soin (diligence), et protéger ce qu'on
apprend (confidentialité).

\subsection{Les principes au quotidien}

\begin{tableaucours}{L{2.8cm}L{4.6cm}L{4.8cm}}{%
    \tableauHeader{Principe} &
    \tableauHeader{Ce qu'il exige} &
    \tableauHeader{Manquement typique}
}
Intégrité & Ne rien affirmer de faux, ne rien taire d'important & Retirer
un constat à la demande du client \\
Objectivité & Juger sur les preuves, sans influence & Adoucir une
conclusion pour conserver une mission \\
Compétence & N'accepter que ce qu'on sait faire, ou s'entourer & Auditer
un système central bancaire sans en connaître l'architecture \\
Diligence & Respecter les normes et le plan, documenter & Conclure sur un
échantillon trop petit pour finir à temps \\
Confidentialité & Ne divulguer ni utiliser l'information hors mission &
Évoquer les failles d'un client devant un autre \\
\end{tableaucours}
\vspace{-0.8em}
\begin{center}{\small\itshape Tableau~16.2. Les principes au
quotidien}\end{center}

\subsection{Pourquoi la déontologie est une condition de la preuve}

L'auditologie pose que les conclusions d'un audit doivent pouvoir être
défendues devant un tiers (axiome~6, opposabilité), et que l'audité et
l'auditeur poursuivent des intérêts distincts (axiome~4, interaction
stratégique)~\cite{minka2026auditologie}. Un auditeur qui dépend de
l'audité, ou qui tait un fait sous pression, produit des conclusions qu'un
tiers peut contester sur ce seul fondement~: la déontologie n'est pas une
morale ajoutée à la technique, c'est une condition de la valeur de la
preuve.

% ============================================================================
\section{L'Indépendance et ses Menaces}
\label{sec:c16-menaces}
% ============================================================================

\subsection{Indépendance de fait, indépendance apparente}

\begin{boxdef}
\termecle{Indépendance de fait}\index{indépendance}~: état d'esprit qui
permet de formuler une conclusion sans être influencé.

\medskip
\termecle{Indépendance apparente}~: absence de circonstances qui
conduiraient un tiers informé à douter de l'indépendance de fait.
\end{boxdef}

La seconde compte autant que la première~: un auditeur parfaitement
honnête, mais qui audite le système qu'il a conçu, ne convaincra pas un
tiers. C'est l'apparence qui fonde l'opposabilité.

\subsection{Cinq menaces, des sauvegardes}

Le code de l'IESBA classe les menaces en cinq catégories~\cite{iesba2018code}~;
le Chapitre~\ref{chap:normes} les a présentées. Le Tableau~16.3 les relie
aux sauvegardes possibles.

\begin{tableaucours}{L{2.6cm}L{4.2cm}L{5.4cm}}{%
    \tableauHeader{Menace} &
    \tableauHeader{Exemple} &
    \tableauHeader{Sauvegarde}
}
Intérêt personnel & Honoraires futurs promis par M.~ETEME & Refus
explicite~; information de l'associé~; plafonnement de la part d'un client
dans le chiffre d'affaires \\
Autorévision & AUDIX CONSEIL audite les profils qu'il a conçus & Confier
l'audit à une autre équipe ou un autre cabinet \\
Représentation & L'auditeur défend la position du client devant la COBAC &
Refuser le rôle de défenseur pendant une mission d'assurance \\
Familiarité & Même chef de mission depuis huit ans chez le même client &
Rotation des responsables de mission \\
Intimidation & Menace de rompre le contrat si le constat reste & Escalade
au comité d'audit~; trace écrite \\
\end{tableaucours}
\vspace{-0.8em}
\begin{center}{\small\itshape Tableau~16.3. Menaces sur l'indépendance et
sauvegardes}\end{center}

\begin{boxalert}
\textbf{Violation typique.} Considérer qu'une menace n'existe pas parce que
l'auditeur se sait honnête.

\smallskip
\textbf{Règle~:} la menace s'apprécie du point de vue d'un tiers informé.
Elle est soit éliminée, soit ramenée à un niveau acceptable par une
sauvegarde écrite~; à défaut, l'auditeur refuse ou quitte la mission.
\end{boxalert}

\subsection{La lecture du cours}

Le Chapitre~\ref{chap:normes} en a tiré une règle de cotation~: un \OUI{}
donné par un auditeur en situation de conflit se cote \NEK{} jusqu'à la
revue par un tiers. L'indépendance n'est donc pas seulement une question
morale~: son défaut se voit dans les mesures, par une hausse de $U$.

% ============================================================================
\section{Résoudre un Dilemme}
\label{sec:c16-dilemme}
% ============================================================================

\subsection{Une méthode en six étapes}

Un dilemme déontologique ne se résout pas par l'intuition, sous pression,
au téléphone. On le résout par écrit.

\begin{boxprocedure}
\textbf{Résoudre un dilemme déontologique}
\begin{enumerate}
    \item \textbf{Établir les faits}~: qui demande quoi, par quel moyen, à
          quelle date~? (Une demande orale est consignée dans une note.)
    \item \textbf{Identifier les principes en jeu} (Tableau~16.2) et les
          menaces (Tableau~16.3).
    \item \textbf{Identifier les parties}~: le client, les tiers qui se
          fieront au rapport, le public.
    \item \textbf{Lister les options}, y compris refuser et quitter la
          mission.
    \item \textbf{Consulter}~: l'associé, le responsable de la qualité du
          cabinet, et si besoin un avocat ou l'ordre professionnel.
    \item \textbf{Décider et documenter}~: la décision, ses motifs, les
          personnes consultées, classées au dossier.
\end{enumerate}
\end{boxprocedure}

\subsection{Le dilemme de M.~BEYALA}

\begin{boxscenario}{Opération LIZBIZ. Retirer deux constats~?}
\textbf{Faits.} Appel du DG le 12~juin~: demande de retrait de deux
observations majeures, et promesse de missions futures. M.~BEYALA consigne
l'appel dans une note le jour même.

\medskip
\textbf{Principes et menaces.} Intégrité (ne pas taire un fait
important)~; objectivité~; menaces d'intérêt personnel (missions promises)
et d'intimidation implicite.

\medskip
\textbf{Parties.} Le DG~; le conseil d'administration, commanditaire
ultime~; le commissaire aux comptes~; les investisseurs de la levée de
fonds, qui se fieront aux comptes.

\medskip
\textbf{Options.} (a) Retirer les constats~: exclu (intégrité). (b) Les
maintenir et transmettre le rapport au seul DG~: insuffisant si la lettre de
mission désigne le conseil. (c) Les maintenir, rappeler au DG la lettre de
mission et la conduite à tenir devant des indices de fraude
(Chapitre~\ref{chap:fraude}), et adresser le rapport au destinataire prévu.
(d) Quitter la mission~: seulement si (c) est rendu impossible.

\medskip
\textbf{Décision.} Option (c). Courrier au DG, copie à l'associée du
cabinet~; refus écrit de toute discussion sur les missions futures tant que
la mission en cours n'est pas close.
\end{boxscenario}

\begin{boxmetier}{Auditeur SI}
\textbf{Refuser sans rompre}
\begin{enumerate}
    \item Répondre par écrit, sur un ton factuel, sans jugement sur
          l'intention du demandeur.
    \item Rappeler la lettre de mission (destinataires, périmètre) et les
          normes appliquées.
    \item Proposer ce qui est possible~: un délai pour la réponse de
          l'audité, une présentation au conseil avec la direction.
    \item Ne jamais négocier un constat contre un avantage.
\end{enumerate}
\end{boxmetier}

% ============================================================================
\section{La Responsabilité de l'Auditeur}
\label{sec:c16-responsabilite}
% ============================================================================

\subsection{Trois responsabilités}

\begin{tableaucours}{L{2.6cm}L{4.4cm}L{5.2cm}}{%
    \tableauHeader{Responsabilité} &
    \tableauHeader{Fondement} &
    \tableauHeader{Exemple}
}
Civile & Faute dans l'exécution du contrat ou envers un tiers, causant un
préjudice & Un investisseur perd de l'argent sur la foi d'un rapport qui a
omis un constat connu \\
Pénale & Infraction (complicité, faux, violation du secret
professionnel) & Un auditeur divulgue les failles d'un client à un
concurrent \\
Disciplinaire & Manquement au code d'un ordre ou d'une association
professionnelle & Retrait de la certification CISA pour violation du code
de l'ISACA \\
\end{tableaucours}
\vspace{-0.8em}
\begin{center}{\small\itshape Tableau~16.4. Les trois responsabilités de
l'auditeur (principes généraux~; les textes applicables au Cameroun sont à
vérifier)}\end{center}

\subsection{Une obligation de moyens}

L'auditeur ne garantit pas l'absence d'anomalie~: il s'engage à mettre en
œuvre les diligences que la norme exige (Chapitre~\ref{chap:droit}). Sa
meilleure défense est son dossier~: un papier de travail qui montre que le
test a été conçu, exécuté et revu selon la norme protège l'auditeur, même
si une fraude n'a pas été vue. À l'inverse, un \OUI{} sans preuve est
indéfendable.

\begin{boxinfo}
\textbf{L'expert judiciaire.} Un auditeur des SI peut être désigné par un
juge comme expert. Son rôle change alors~: il répond aux questions posées
par le juge, dans le respect du contradictoire judiciaire, et son rapport
devient une pièce de procédure. Les règles de la preuve numérique
(Chapitre~\ref{chap:audit-preuve}) et celles de l'expertise, propres à
chaque juridiction, s'appliquent (textes camerounais à vérifier).
\end{boxinfo}

% ============================================================================
\section{Certifications et Compétence}
\label{sec:c16-certif}
% ============================================================================

\subsection{Les certifications du métier}

\begin{tableaucours}{L{2.6cm}L{2.6cm}L{6.8cm}}{%
    \tableauHeader{Certification} &
    \tableauHeader{Organisme} &
    \tableauHeader{Domaine}
}
CISA & ISACA & Audit, contrôle et assurance des SI \\
CISM & ISACA & Management de la sécurité de l'information \\
CRISC & ISACA & Risques informatiques et contrôles \\
CGEIT & ISACA & Gouvernance des SI \\
CDPSE & ISACA & Protection des données dès la conception \\
CIA & IIA & Audit interne \\
Auditeur principal ISO/IEC~27001 & Organismes de formation accrédités &
Audit de SMSI selon l'ISO~19011 et l'ISO/IEC~27007 \\
\end{tableaucours}
\vspace{-0.8em}
\begin{center}{\small\itshape Tableau~16.5. Les principales certifications
de l'audit et de la sécurité des SI}\end{center}

Une certification n'est ni un diplôme ni une licence d'exercer~: elle
atteste un niveau de connaissance à une date, et engage son titulaire à
respecter un code et à se former. L'ISACA exige ainsi de ses certifiés un
volume minimal annuel et triennal d'heures de formation professionnelle
continue~\cite{isaca2024ethique}.

\subsection{La compétence, principe déontologique}

Accepter une mission qu'on ne sait pas mener est un manquement, pas une
simple erreur. Le chef de mission vérifie, avant d'accepter, que l'équipe
réunit les compétences nécessaires (bases de données, réseau, système
central bancaire, droit applicable)~; à défaut, il s'adjoint un expert,
dont il reste responsable (Chapitre~\ref{chap:mission}).

\begin{boxmetier}{Responsable de la conformité}
\textbf{Vérifier les qualifications des prestataires d'audit.} Lorsque la
BANQUE SANAGA confie un audit à un cabinet, le responsable de la
conformité vérifie les certifications annoncées (registre en ligne des
organismes), la déclaration d'indépendance de chaque membre de l'équipe,
et l'absence de mission de conseil récente sur le même objet.
\end{boxmetier}

% ============================================================================
\section{Le Marché de l'Audit}
\label{sec:c16-marche}
% ============================================================================

\subsection{Qui audite}

Au Cameroun comme ailleurs, l'audit des SI est exercé par~: les services
d'audit interne des grandes organisations (banques, opérateurs, entreprises
publiques)~; les cabinets d'audit et d'expertise comptable, dont les
réseaux internationaux~; des cabinets spécialisés en sécurité~; les
organes publics de contrôle (Chambre des comptes, organes de contrôle
supérieur de l'État)~; les autorités de régulation et leurs auditeurs
agréés. L'exercice du commissariat aux comptes est réservé aux
professionnels inscrits à l'ordre (ONECCA au Cameroun~; conditions à
vérifier).

\subsection{Les tensions du marché}

\begin{itemize}
    \item \textbf{Le client paie l'auditeur}~: c'est la source permanente de
          la menace d'intérêt personnel. Les sauvegardes~: comité d'audit
          qui choisit et rémunère, plafonnement de la dépendance
          économique, rotation.
    \item \textbf{Audit et conseil chez le même client}~: la vente de
          conseil après l'audit crée l'autorévision. L'affaire Enron
          (Chapitre~\ref{chap:defaillances}) a conduit la loi
          Sarbanes-Oxley à restreindre ces cumuls pour les auditeurs légaux
          des sociétés cotées américaines.
    \item \textbf{La pression sur les prix}~: un audit vendu trop bas se
          paie en diligence. Le chef de mission doit refuser un budget qui
          ne permet pas de respecter la norme, ou réduire le périmètre par
          écrit.
\end{itemize}

\begin{boxscenario}{Opération SANAGA. L'appel d'offres}
\textbf{Faits.} Pour son audit de sécurité annuel, la BANQUE SANAGA reçoit
trois offres~: 18, 32 et 35~millions FCFA. L'offre la moins chère prévoit
6~jours de terrain pour un système central, 40~agences et un prestataire.

\medskip
\textbf{Analyse} (confiance moyenne). Six jours ne permettent ni de tester
les contrôles généraux du système central, ni d'examiner le prestataire~:
l'offre ne peut produire que des déclarations. Le comité d'audit retient
l'offre de 32~millions, et le motive par écrit~: la diligence a un coût.
\end{boxscenario}

\begin{boxmetier}{Gestionnaire de risques}
\textbf{Traiter la dépendance à un auditeur comme un risque.} Une banque
qui confie depuis dix ans toutes ses missions d'audit et de conseil au
même cabinet porte un risque de familiarité et d'autorévision. Le
gestionnaire de risques l'inscrit au registre et propose une rotation.
\end{boxmetier}

\separateur

\begin{boxresume}
\textbf{Ce qu'il faut retenir de ce chapitre~:}
\begin{itemize}
    \item \textbf{Trois codes} (IIA, ISACA, IESBA) convergent sur
          l'intégrité, l'objectivité, la compétence, la diligence et la
          confidentialité.
    \item \textbf{L'indépendance} est de fait et d'apparence~; cinq menaces
          (intérêt personnel, autorévision, représentation, familiarité,
          intimidation) appellent des sauvegardes écrites, ou le refus.
    \item \textbf{La déontologie est une condition de la preuve}~: un
          auditeur dépendant produit des conclusions contestables~; un
          \OUI{} en conflit se cote \NEK{}.
    \item \textbf{Un dilemme se résout par écrit}, en six étapes, avec
          consultation.
    \item \textbf{L'auditeur engage sa responsabilité} civile, pénale et
          disciplinaire~; son dossier est sa défense.
    \item \textbf{Les certifications} attestent une compétence et engagent à
          un code et à la formation continue.
    \item \textbf{Le marché de l'audit} crée des tensions (client payeur,
          cumul audit et conseil, prix) que des sauvegardes organisées
          compensent.
\end{itemize}
\end{boxresume}

\begin{boxinfo}
\textbf{Points de contrôle de l'élève.} Je sais~: citer les cinq exigences
communes aux codes~; nommer et illustrer les cinq menaces~; proposer une
sauvegarde~; appliquer la méthode en six étapes~; distinguer les trois
responsabilités~; citer les certifications de l'ISACA et de l'IIA.
\end{boxinfo}

\bigskip

\begin{boxexo}[title={\ding{45}\quad Exercice~16.1. Quel principe~? \niveaubase}]
Indiquez le principe déontologique en jeu~: (a) un auditeur accepte
l'audit d'un réseau industriel sans jamais en avoir vu~; (b) une auditrice
garde sur sa clé personnelle les extractions d'un client~; (c) un chef de
mission conclut sur 10 paiements au lieu des 60 prévus pour tenir les
délais~; (d) un rapport passe sous silence un constat sur un administrateur
de la société.
\end{boxexo}

\begin{boxexo}[title={\ding{45}\quad Exercice~16.2. Quelle menace~? \niveaubase}]
Identifiez la menace et proposez une sauvegarde~: (a) la fille du chef de
mission travaille à la DSI auditée~; (b) l'auditeur a rédigé la politique
de sécurité qu'il audite~; (c) le client menace de faire réaliser
l'audit suivant par un concurrent~; (d) le cabinet tire 40~\% de son
chiffre d'affaires d'un seul client.
\end{boxexo}

\begin{boxexo}[title={\ding{45}\quad Exercice~16.3. Fait ou apparence \niveaumoyen}]
Mme~EKAMBI a travaillé à la DSI de la BANQUE SANAGA jusqu'en 2024. Elle
affirme n'avoir eu aucun rôle dans la migration de 2026. Son indépendance
de fait est-elle menacée~? son indépendance apparente~? Quelle sauvegarde
proposez-vous pour l'audit des changements~?
\end{boxexo}

\begin{boxexo}[title={\ding{45}\quad Exercice~16.4. Les six étapes \niveaumoyen}]
Appliquez la méthode en six étapes au cas suivant~: le magistrat
vérificateur M.~NJOYA reçoit, de la part d'un conseiller du ministre, une
demande « amicale » de « ne pas insister » sur les paiements après départ
de l'ANPP.
\end{boxexo}

\begin{boxexo}[title={\ding{45}\quad Exercice~16.5. Le courrier de M.~BEYALA \niveaumoyen}]
Rédigez le courrier de M.~BEYALA au directeur général de LiZbiZ, suite à
l'appel du 12~juin, en respectant la méthode « refuser sans rompre ».
\end{boxexo}

\begin{boxexo}[title={\ding{45}\quad Exercice~16.6. La cotation en conflit \niveaumoyen}]
Dans l'audit de suivi de LiZbiZ, AUDIX CONSEIL doit coter six contrôles
portant sur les profils de droits qu'il a lui-même conçus. Il les jugerait
tous \OUI{}. Comment les cotez-vous dans le rapport, et quel est l'effet
sur $U$ du domaine (20 réponses, 18/1/1 avant la prise en compte du
conflit)~?
\end{boxexo}

\begin{boxexo}[title={\ding{45}\quad Exercice~16.7. Responsabilité \niveaumoyen}]
Un investisseur poursuit AUDIX CONSEIL après la découverte de la fraude
de TRANSLOG SERVICES. Que doit contenir le dossier du cabinet pour établir
qu'il a respecté son obligation de moyens~? Dans quel cas sa
responsabilité serait-elle difficile à écarter~?
\end{boxexo}

\begin{boxexo}[title={\ding{45}\quad Exercice~16.8. Une certification \niveaubase}]
Un élève ingénieur veut devenir auditeur des SI dans une banque.
Quelle certification vise-t-il en premier, et pourquoi~? Que lui impose-t-
elle une fois obtenue~?
\end{boxexo}

\begin{boxexo}[title={\ding{45}\quad Exercice~16.9. Audit et conseil \niveauexpert}]
La BANQUE SANAGA veut confier au même cabinet l'audit de sécurité annuel,
la conception de son plan de continuité et l'accompagnement de sa
certification ISO/IEC~27001. Analysez les menaces, proposez une
organisation acceptable, et dites ce que le comité d'audit doit décider et
documenter.
\end{boxexo}

\begin{boxexo}[title={\ding{45}\quad Exercice~16.10. Un code pour AUDIX CONSEIL \niveauexpert}]
Rédigez, en dix règles au plus, le code de conduite interne d'AUDIX
CONSEIL pour ses missions d'audit des SI~: acceptation des missions,
indépendance, cumul audit et conseil, cadeaux et invitations, données des
clients, réseaux sociaux, signalement d'un manquement d'un collègue.
\end{boxexo}

\bigskip

\begin{boxinfo}
\textbf{Pour aller plus loin.}
\begin{itemize}
    \item IIA, Normes mondiales d'audit interne, domaine~II~\cite{iia2024gias}.
    \item ISACA, \textit{Code of Professional Ethics}~\cite{isaca2024ethique}.
    \item IESBA, \textit{International Code of Ethics for Professional
          Accountants}~\cite{iesba2018code}~: les menaces et les
          sauvegardes, en détail.
\end{itemize}
\end{boxinfo}

\begin{boxinfo}
\textbf{Transition.} La Partie~\ref{part:normatif} a fixé les normes, les
référentiels, le droit et la déontologie. La Partie~\ref{part:pratique}
applique la méthode aux grands domaines du système d'information~: la
gouvernance et la DSI d'abord (Chapitre~\ref{chap:audit-gouvernance}).
\end{boxinfo}
